%% kapa: exam=1 round=30 year=2019 period=1 form=A subject=real_estate subject_ko=부동산학원론 seq=4 number=84
%% source: Q-Net 감정평가사 기출문제 (cst00309) artlSeq=5205205 file=extracted/1008/5205205/2019년 제30회 감정평가사 1차시험 문제지/2019년 제30회 감정평가사 1차시험 문제지/2019년도 제30회 감정평가사 1차 1교시 A형.hwp
%% conversion: kordoc 4.24.51 -> tools/gen1.py (auto); review=proofread-sonnet; flags=none
\begin{question}{84}
부동산의 개념에 관한 설명으로 옳지 않은 것은?
\begin{choices}
  \item 토지는 제품생산에 필요한 부지를 제공하는 생산요소이다.
  \item 토지는 생활의 편의를 제공하는 최종 소비재이기도 하다.
  \item ｢민법｣상 부동산은 토지 및 그 정착물이며, 부동산 이외의 물건은 동산이다.
  \item 준부동산에는 등기나 등록수단으로 공시된 광업재단, 공장재단, 선박, 항공기, 어업권 등이 있다.
  \item ｢입목에 관한 법률｣에 의해 소유권보존등기를 한 입목은 토지와 분리하여 양도할 수 없다.
\end{choices}
\end{question}
\answer{5}
